% ============================================================================
%  part2a_inference.tex  --  Part II, sections 4-5.
%  Sparse variational Gaussian-process inference + the eigenspace representation.
%  Self-contained \input-ready fragment; uses ONLY the shared preamble macros.
% ============================================================================

\section{Sparse variational Gaussian-process inference}\label{sec:varinf}

Part~I fixed the generative model: a zero-mean latent Gaussian process
$\lat \sim \GProc(0,\Kf)$ over input space, an exponential link to the rate
$\rate(x)=\exp\!\big(\gain\,\lat(x)+\latz\big)$, and Poisson counts
$\spk_i \sim \Poi\!\big(\rate(x_i)\big)$, where $\Kf$ is the first-order
arc-cosine kernel of \eqref{eq:acoskernel} built on the structured spatial
prior covariance $\Cmat$ of \eqref{eq:Cprior}. Here we turn that model into a
trainable posterior. The variational posterior developed below can be
represented equivalently in the full inducing space or in the eigenbasis of
\S\ref{sec:eigenspace}.

\subsection{Poisson breaks conjugacy}\label{sec:conjugacy}

Given a batch of inputs $X=(x_1,\dots,x_{\Ntr})$ and counts
$\spk=(\spk_1,\dots,\spk_{\Ntr})$, the exact posterior over the latent and the
marginal likelihood,
\begin{align}
  p(\lat\mid\spk,X,\theta)&\;\propto\;
    \Big[\textstyle\prod_i \Poi\!\big(\spk_i\mid e^{\gain\lat_i+\latz}\big)\Big]\,
    \Gauss(\lat\mid 0,\Kf_\theta), \notag\\
  p(\spk\mid X,\theta)&=\!\int\! \prod_i \Poi\!\big(\spk_i\mid e^{\gain\lat_i+\latz}\big)\,
    \Gauss(\lat\mid 0,\Kf_\theta)\,d\lat ,
\end{align}
have no closed form: the exponential-Poisson likelihood is not conjugate to the
Gaussian prior, so the integral over $\lat$ is intractable. We therefore
approximate the posterior by a tractable Gaussian and maximise a variational
lower bound (\S\ref{sec:elbo}).

\begin{remark}[The conjugate contrast]
For a \emph{Gaussian} observation $\spk=\lat+\varepsilon$,
$\varepsilon\sim\Gauss(0,\pvar_r I)$, everything would be closed-form: the
marginal is $\spk\mid X\sim\Gauss(0,\Kf+\pvar_r I)$, the predictive at a test
input $x_\ast$ is Gaussian with mean $\kvec_\ast\transpose(\Kf+\pvar_r I)^{-1}\spk$
and variance $\Kf(x_\ast,x_\ast)-\kvec_\ast\transpose(\Kf+\pvar_r I)^{-1}\kvec_\ast$,
and the log-evidence is
$-\tfrac12\spk\transpose(\Kf+\pvar_r I)^{-1}\spk-\tfrac12\log\lvert\Kf+\pvar_r I\rvert-\tfrac{\Ntr}{2}\log 2\pi$.
None of these survive the Poisson link; the variational route below recovers a
usable posterior in its place.
\end{remark}

\subsection{Inducing points and the variational posterior}\label{sec:inducing}

\begin{definition}[Sparse variational posterior]
Introduce $\Mind$ \emph{inducing points} $\Ipt=\{\ipt_1,\dots,\ipt_{\Mind}\}$
with latent values $\latt=\lat(\Ipt)$, and place a
Gaussian variational posterior on them,
\begin{equation}
  q(\latt)=\Gauss(\latt\mid\vm,\vV),\qquad
  \vm\in\Reals^{\Mind},\quad \vV\in\mathbb{S}_{++}^{\Mind}.
  \label{eq:qdef}
\end{equation}
The posterior at \emph{any} input $x$ is recovered by marginalising the
conditional prior $p(\lat(x)\mid\latt)$ against $q(\latt)$.
\end{definition}

Exact inference would need the $\Ntr\!\times\!\Ntr$ kernel over all training
inputs, an $O(\Ntr^3)$ cost; the inducing construction reduces it to
$O(\Ntr\Mind^2)$.

\begin{remark}[Active-learning invariant $\Mind=\Ntr$]\label{rem:MeqN}
In the active-learning setting the inducing set is taken equal to the training
set, $X=\Ipt$, so $\Mind=\Ntr$ and $\Kf=\Ktil$. The ``sparse'' posterior is then
\emph{exact}. The SVGP formalism is retained not for sparsity but for two other
reasons: (i)~it supplies the closed-form variational machinery of
\S\ref{sec:predmoments}--\S\ref{sec:fbar}, and (ii)~its inducing structure admits
an $O(\Mind)$ rank-one online growth that appends one input per active-learning
step.
\end{remark}

\subsection{Projected posterior moments}\label{sec:predmoments}

Marginalising $p(\lat(x)\mid\latt)\,q(\latt)$ yields a Gaussian
$q(\lat(x))=\Gauss\!\big(\pmean(x),\pvar(x)\big)$ whose moments are the
inducing-point moments projected through the kernel:
\begin{align}
  \pmean(x)  &= \kvec(x)\transpose\,\Ktil^{-1}\,\vm , \notag\\
  \pvar(x)   &= \Kf(x,x) + \kvec(x)\transpose\,\Ktil^{-1}\big(\vV-\Ktil\big)\Ktil^{-1}\,\kvec(x),
  \label{eq:predmoments}
\end{align}
with the cross-covariance vector $\kvec(x)=\Kf(\Ipt,x)\in\Reals^{\Mind}$, the
inducing gram $\Ktil=\Kf(\Ipt,\Ipt)\in\Reals^{\Mind\times\Mind}$, and the prior
self-variance $\Kf(x,x)$ given by the self-kernel identity \eqref{eq:selfkernel}.
Introducing the \emph{projection vector} $\pv(x)=\Ktil^{-1}\kvec(x)$ and the
Schur complement $s(x)=\Kf(x,x)-\kvec(x)\transpose\Ktil^{-1}\kvec(x)$, the
variance splits into a data-independent residual and an epistemic term,
\begin{equation}
  \pvar(x)=\underbrace{\big[\Kf(x,x)-\kvec(x)\transpose\Ktil^{-1}\kvec(x)\big]}_{\text{residual (Nystr\"om/Schur) }s(x)}
    \;+\;\underbrace{\kvec(x)\transpose\Ktil^{-1}\vV\Ktil^{-1}\kvec(x)}_{\text{epistemic}}
    \;=\; s(x)+\pv(x)\transpose\vV\,\pv(x),
  \label{eq:varsplit}
\end{equation}
and the mean reads $\pmean(x)=\pv(x)\transpose\vm$. The two forms of $\pvar$ are
algebraically identical since
$\kvec\transpose\Ktil^{-1}(\vV-\Ktil)\Ktil^{-1}\kvec=\kvec\transpose\Ktil^{-1}\vV\Ktil^{-1}\kvec-\kvec\transpose\Ktil^{-1}\kvec$.

\begin{remark}[Prior sanity check]
At the prior $\vV=\Ktil$ the correction in \eqref{eq:predmoments} vanishes and
$\pvar(x)=\Kf(x,x)$ (the prior variance); as data drive $\vV$ below $\Ktil$,
$\pvar(x)<\Kf(x,x)$.
\end{remark}

These moments can be evaluated directly in the $\Mind$-dimensional inducing
space or, equivalently, in the eigenbasis of $\Ktil$ (\S\ref{sec:eig-moments}).

\subsection{The ELBO}\label{sec:elbo}

All parameters are learned by maximising the evidence lower bound on
$\log p(\spk)$,
\begin{equation}
  \ELBO \;=\; \sum_i \E_{q(\lat_i)}\!\big[\log \Poi(\spk_i\mid\lat_i)\big]
    \;-\; \KLdiv\!\big(q(\latt)\,\big\|\,p(\latt)\big),
  \label{eq:elbo}
\end{equation}
which is tight when $q$ equals the true posterior. Because the likelihood
factorises over inputs and $q(\lat_i)=\Gauss(\pmean_i,\pvar_i)$ is Gaussian, each
expected log-likelihood term is closed-form:
\begin{equation}
  \E_{q(\lat_i)}\!\big[\log \Poi(\spk_i\mid\lat_i)\big]
    \;=\; \spk_i\big(\gain\pmean_i+\latz\big)
      \;-\; \exp\!\big(\gain\pmean_i+\tfrac12\gain^2\pvar_i+\latz\big)
      \;+\; \text{const},
  \label{eq:ell}
\end{equation}
with $\text{const}=-\log(\spk_i!)$ independent of all parameters. The KL term is
the standard Gaussian--Gaussian divergence between $q(\latt)=\Gauss(\vm,\vV)$ and
the prior $p(\latt)=\Gauss(0,\Ktil)$,
\begin{equation}
  \KLdiv(q\,\|\,p) \;=\; \tfrac12\Big[\,
    \log\frac{\lvert\Ktil\rvert}{\lvert\vV\rvert}
    \;+\; \Tr\!\big(\Ktil^{-1}\vV\big)
    \;+\; \vm\transpose\Ktil^{-1}\vm
    \;-\; \Mind \,\Big].
  \label{eq:kl}
\end{equation}

\begin{remark}[The $-\Mind$ constant]
The $-\Mind$ in \eqref{eq:kl} is parameter-independent (it is the dimension of
the inducing Gaussian) and carries zero gradient, so it affects only the reported
loss magnitude, never the optimum or the predictions. The eigenspace evaluation
of \S\ref{sec:eig-kl} carries the analogous $-\nb$ for the same reason.
\end{remark}

The full objective is maximised by EM-style coordinate ascent
(E-step on $(\vm,\vV)$, F-step on $(\gain,\latz)$, M-step on the kernel
hyperparameters); the update equations are derived in the training-algorithm
section. Equation~\eqref{eq:elbo} is the single objective all three blocks share.

\subsection{Expected rate via the Gaussian MGF}\label{sec:fbar}

The exponential term in \eqref{eq:ell} is the \emph{expected rate} under $q$,
\begin{equation}
  \fbar_i \;=\; \E_{q(\lat_i)}\!\big[\exp(\gain\lat_i+\latz)\big]
    \;=\; \exp\!\big(\gain\pmean_i+\tfrac12\gain^2\pvar_i+\latz\big).
  \label{eq:fbar}
\end{equation}
This is the Gaussian moment-generating function
$\E[e^{tZ}]=\exp(t\mu+\tfrac12 t^2\sigma^2)$ applied to $Z=\gain\lat_i+\latz$ at
$t=1$.

\begin{remark}[Jensen inflation]
The term $\tfrac12\gain^2\pvar_i$ is strictly positive whenever $\pvar_i>0$, so
posterior \emph{uncertainty raises} the predicted rate above the point estimate
$\exp(\gain\pmean_i+\latz)$ --- a direct consequence of Jensen's inequality for
the convex map $\exp$. Confidence and rate are thus coupled: a poorly-constrained
input is predicted to have a \emph{higher} rate, not lower.
\end{remark}

\subsection{Prediction: expected count}\label{sec:prediction}

The moments \eqref{eq:predmoments} do not depend on the observed $\spk$, so the
same formulas serve as the predictive distribution at any novel input $x_\ast$.
The expected count is the expected rate \eqref{eq:fbar} evaluated there,
\begin{equation}
  \langle \spk_\ast\rangle=\langle\rate(x_\ast)\rangle
    =\exp\!\big(\gain\pmean(x_\ast)+\tfrac12\gain^2\pvar(x_\ast)+\latz\big),
\end{equation}
with $\pmean(x_\ast),\pvar(x_\ast)$ from \eqref{eq:predmoments} using freshly
evaluated $\kvec(x_\ast)$ and $\Kf(x_\ast,x_\ast)$. The rate is
log-normal, so its variance is
\begin{equation}
  \operatorname{Var}\!\big(\rate(x)\big)
    =\langle\rate(x)\rangle^2\big(e^{\gain^2\pvar(x)}-1\big).
  \label{eq:fvar}
\end{equation}

\subsection{The posterior cross-covariance pitfall}\label{sec:crosscov}

The self-variance of \eqref{eq:predmoments} extends formally to the
\emph{cross-covariance} between two distinct inputs $x$ (a candidate sample
location) and $x_\ast$ (a query location):
\begin{equation}
  \Sigma_{x,x_\ast}=\operatorname{Cov}\!\big[\lat(x),\lat(x_\ast)\mid \mathcal D\big]
    =\Kf(x,x_\ast)+\pv(x)\transpose\big(\vV-\Ktil\big)\pv(x_\ast),
  \qquad \pv(\cdot)=\Ktil^{-1}\kvec(\cdot),
  \label{eq:crosscov}
\end{equation}
which reduces to $\pvar(x)$ at $x=x_\ast$.

\begin{remark}[Failure outside the inducing hull]
For a query $x_\ast$ \emph{outside} the inducing region, \eqref{eq:crosscov}
produces cross-covariances that violate the Cauchy--Schwarz bound
$\lvert\Sigma_{x,x_\ast}\rvert\le\sqrt{\Sigma_{x,x}\,\Sigma_{x_\ast,x_\ast}}$.
Measured example (RBF, $\ell=0.2$, $20$ inducing points in $[-1,1]$, sample
$x=0$): at $x_\ast=\pm1.5$ the manual formula gives $\Sigma\approx-4$ while
the consistent joint covariance gives $\approx 0$; with self-variances $\sim0.04$
the valid maximum is $\sim0.04$, so $-4$ is two orders of magnitude out of range.
Root cause: in the residual $\Kf(x,x_\ast)-\pv(x)\transpose\Ktil\,\pv(x_\ast)$ the
direct kernel $\Kf(x,x_\ast)\!\to\!0$ (finite lengthscale) outside the hull, but
$\pv(x_\ast)=\Ktil^{-1}\kvec(x_\ast)$ can have large entries through $\Ktil^{-1}$,
and the two pieces fail to cancel.
\end{remark}

Left unchecked, a value like $\Sigma_{x,x_\ast}\approx-4$ feeds the conditional
moments
$\pmean_{\mathrm{cond}}(x_\ast)=\pmean(x_\ast)+(\Sigma_{x,x_\ast}/\Sigma_{x,x})(\lat_{\mathrm{obs}}-\pmean(x))$
and $\pvar_{\mathrm{cond}}(x_\ast)=\Sigma_{x_\ast,x_\ast}-\Sigma_{x,x_\ast}^2/\Sigma_{x,x}$
of the distribution-aware utility (Part~III) with a regression coefficient
$\approx-100$, giving nonsensical conditional means and spurious utility peaks at
the edges of the inducing region. The remedy is to never assemble the cross-block
from the manual formula \eqref{eq:crosscov}; instead read the cross-block from the
jointly-formed covariance of $(x,x_\ast)$, which is Nystr\"om-consistent --- it
carries the conditional correction term that \eqref{eq:crosscov} drops outside the
hull. The pathology is therefore confined to hand-assembled cross-covariances; any
computation that forms a cross-covariance from \eqref{eq:crosscov} outside the
inducing hull will reintroduce it.

\subsection{Whitening}\label{sec:whitening}

An alternative reparameterisation is \emph{Cholesky whitening}, which makes the
prior a standard normal. With the Cholesky factor $\Ktil=\Lchol\Lchol\transpose$,
define the whitened inducing values
\begin{equation}
  \widetilde{\latt}=\Lchol^{-1}(\latt-\mu_z)\;\sim\;\Gauss(0,I)\quad\text{under the prior},
\end{equation}
where $\mu_z$ is the prior mean at the inducing points. Optimisation conditions
much better in these coordinates. In these coordinates one stores the
\emph{whitened} variational parameters $(\widetilde{\vm},\widetilde{\vV})$,
related to the actual ones by
\begin{equation}
  \vm-\mu_z=\Lchol\,\widetilde{\vm},\quad \vV=\Lchol\,\widetilde{\vV}\,\Lchol\transpose;
  \qquad
  \widetilde{\vm}=\Lchol^{-1}(\vm-\mu_z),\quad \widetilde{\vV}=\Lchol^{-1}\vV\,\Lchol^{-\mathsf{T}}.
  \label{eq:whiten}
\end{equation}
The whitened mean thus encodes the \emph{deviation} from the prior mean function,
not absolute function values. Recovering the moments from the whitened parameters
therefore requires three corrections: unwhiten the covariance
($\vV=\Lchol\widetilde{\vV}\Lchol\transpose$), unwhiten the mean
($\vm-\mu_z=\Lchol\widetilde{\vm}$), and add the mean function back,
$\pmean_\ast=\pmean(x_\ast)+\pv(x_\ast)\transpose\Lchol\widetilde{\vm}$ and
$\pvar_\ast=s(x_\ast)+\pv(x_\ast)\transpose\Lchol\widetilde{\vV}\Lchol\transpose\pv(x_\ast)$.

\begin{remark}[Zero prior mean]
With a zero prior mean, $\mu_z=0$ and the ``add the mean function'' step is
trivially zero; the whitening maps \eqref{eq:whiten} themselves apply for any
prior mean.
\end{remark}


\section{The eigenspace representation}\label{sec:eigenspace}

The variational posterior $q(\latt)$ of \S\ref{sec:varinf} can be represented
equivalently in the full $\Mind$-dimensional inducing space or, more
economically, in the eigenbasis of the inducing gram $\Ktil$. This section
develops the eigenbasis form: the same prior $\Gauss(0,\Ktil)$, the same
arc-cosine kernel, the same $\Poi(\gain,\latz)$ likelihood, the same expected
rate \eqref{eq:fbar}, and the same ELBO \eqref{eq:elbo}, re-expressed in a
low-dimensional basis in which $\Ktil$ is diagonal.

\subsection{Eigendecomposition of the inducing gram}\label{sec:eigdecomp}

The inducing gram $\Ktil$ ($\Mind\times\Mind$) has effective rank far below
$\Mind$ (typically $\sim10$--$50$). We eigendecompose it and keep only the
well-conditioned directions:
\begin{equation}
\begin{gathered}
  \Ktil=\eb\,\evals\,\eb\transpose,\qquad
  \evals=\diago(\text{eigvals}_b), \\
  \eb\in\Reals^{\Mind\times\nb},\qquad \eb\transpose\eb=I_{\nb},
\end{gathered}
  \label{eq:eigdecomp}
\end{equation}
where the kept set is chosen by a \emph{relative} threshold with an absolute
floor,
\begin{equation}
  \text{keep eigenpair }j \iff \evals_j > \max\!\big(\evals_{\max}\cdot\texttt{tol},\ \texttt{tol}\big),
  \qquad \texttt{tol}=10^{-4}.
  \label{eq:eigthresh}
\end{equation}
The eigenvalues are taken in ascending order; $\eb$ collects the kept eigenvectors
as columns and $\evals$ the kept eigenvalues. The kept count $\nb$ is
\emph{dynamic}: it shifts across training iterations as the kernel hyperparameters
move, and grows by $\sim1$ when a point is appended. Typical $\nb\approx10$--$11$
in the active-learning setting, up to $\sim50$ in larger problems.

Three payoffs motivate the projection: (i)~dimensionality drops $\Mind\to\nb$, so
per-step cost falls $O(\Mind^3)\to O(\nb^3)$; (ii)~$\Ktil$ becomes
\emph{diagonal} in the new basis, $\Ktilb=\evals$, so $\Ktil^{-1}$ is the
element-wise reciprocal $\diago(1/\evals)$ --- no matrix solve; (iii)~dropping the
small eigenvalues gives implicit regularisation.

\subsection{The eigenspace state}\label{sec:eigstate}

The eigenspace posterior is carried together with the precomputed kernel
projections used below:

\begin{center}\small
\begin{tabular}{@{}l l >{\raggedright\arraybackslash}p{8.6cm}@{}}
\toprule
symbol & shape & meaning \\
\midrule
$\mb$ & $(\nb)$ & variational mean in the eigenbasis, $\mb=\eb\transpose\vm$ \\
$\Vb$ & $(\nb,\nb)$ & variational covariance in the eigenbasis, $\Vb=\eb\transpose\vV\eb$ --- \textbf{dense} \\
$\eb$ & $(\Mind,\nb)$ & eigenvector matrix of $\Ktil$ (the basis) \\
$\evals$ & $(\nb)$ & kept eigenvalues of $\Ktil$ (ascending) \\
$\Ktilb$ & $(\nb,\nb)$ & inducing gram in the eigenbasis $=\diago(\evals)$ --- \textbf{diagonal} \\
$\Kf_b$ & $(\Ntr,\nb)$ & cross-kernel projected into the eigenbasis, $\Kf_b=\Kf(X,\Ipt)\,\eb$ \\
$\amat$ & $(\Ntr,\nb)$ & projection $\Kf\Ktil^{-1}$ in the eigenbasis, $\amat=\Kf_b/\evals$ (element-wise) \\
$\Kvec$ & $(\Ntr)$ & diagonal self-kernel $\Kf(x_i,x_i)$, the per-point prior variance \\
\bottomrule
\end{tabular}
\end{center}

\begin{remark}[Only $\Ktilb$ is diagonal]
$\Vb$ is a \emph{dense} $\nb\times\nb$ matrix even though $\Ktilb$ is diagonal;
never assume a diagonal $\Vb$. The projection row $\amat_i=(\Kf_b/\evals)_i$
is exactly $\kvec(x_i)\transpose\Ktil^{-1}$ expressed in the eigenbasis, obtained
by element-wise division rather than a matrix solve.
\end{remark}

Initialisation sets the variational parameters to the prior:
\begin{equation}
  \mb=0,\qquad \Vb=\Ktilb \quad(\text{i.e.\ }\vm=0,\ \vV=\Ktil,\ \text{so } \KLdiv=0\text{ at start}).
\end{equation}

\subsection{Posterior moments in the eigenbasis}\label{sec:eig-moments}

With $\amat=\Kf_b/\evals$ (element-wise), the projected moments
\eqref{eq:predmoments} at the training points become:
\begin{align}
  \pmean &= \amat\,\mb, \notag\\
  \pvar  &= \Kvec + \diago\!\big(\amat\,(\Vb-\Ktilb)\,\amat\transpose\big)
          = \Kvec + \operatorname{rowsum}\!\big(\amat\odot(\amat\,(\Vb-\Ktilb))\big), \notag\\
  \pvar  &\leftarrow \max\!\big(\pvar,\ 10^{-6}\big),
  \label{eq:eigmoments}
\end{align}
where $\pmean,\pvar$ are the per-training-point posterior-moment vectors and
$\odot$ is the element-wise product; the two variance expressions coincide
because $\operatorname{rowsum}(\amat\odot(\amat M))=\diago(\amat M\amat\transpose)$
for any $M$. Term by term this equals \eqref{eq:predmoments}: $\Kvec_i=\Kf(x_i,x_i)$
(the self-kernel identity \eqref{eq:selfkernel}) and
$\amat_i(\Vb-\Ktilb)\amat_i\transpose=\kvec(x_i)\transpose\Ktil^{-1}(\vV-\Ktil)\Ktil^{-1}\kvec(x_i)$.
For \textbf{test} points the same two expressions run on a freshly projected
$\amat=\Kf(x_\ast,\Ipt)\eb/\evals$ with $\Kvec=\Kf(x_\ast,x_\ast)$, so no
training-point kernel is recomputed; the predicted rate there is
$\rate_{\mathrm{pred}}=\exp(\gain\pmean+\tfrac12\gain^2\pvar+\latz)$, i.e.\
\eqref{eq:fbar}.

\subsection{ELBO and KL in the eigenbasis}\label{sec:eig-kl}

The eigenspace ELBO uses the log-likelihood \eqref{eq:ell} and, because
$\Ktilb=\diago(\evals)$, the KL \eqref{eq:kl} collapses to sums over the kept
dimensions:
\begin{align}
  \KLdiv(q\,\|\,p)
    &= \tfrac12\Big[\Tr(\Ktil^{-1}\vV)+\vm\transpose\Ktil^{-1}\vm-\nb
        +\log\lvert\Ktil\rvert-\log\lvert\vV\rvert\Big] \notag\\
    &= \tfrac12\Big[\ \textstyle\sum_j \dfrac{\Vb[j,j]}{\evals_j}
        \;+\; \sum_j \dfrac{\mb[j]^2}{\evals_j}
        \;-\; \nb
        \;+\; \sum_j \log\evals_j
        \;-\; \log\det\Vb\ \Big].
  \label{eq:eigkl}
\end{align}

\begin{remark}[Diagonal-trace assumption]
The trace term uses only $\diago(\Vb)$, which is valid \emph{only because}
$\Ktilb$ is diagonal in the basis $\eb$ in which the ELBO is evaluated --- i.e.\
when the kernel and the basis are mutually consistent. The $\log\lvert\vV\rvert$
term uses the full log-determinant of the dense $\Vb$. The $-\nb$ is a
zero-gradient constant (cf.\ the $-\Mind$ remark after \eqref{eq:kl}): it shifts
only the reported loss magnitude.
\end{remark}
